\section{Further results and choice implications}\label{app:auxiliary-results}

This appendix follows the order of Sections~\ref{sec:theorem}
and~\ref{sec:implications}.  It first proves the representation's immediate
consequences and the independence result, then turns to the trace test,
deliberate randomisation, record value, and measurement.  None of these
results is used to prove Theorem~\ref{thm:main}; Appendix~\ref{app:main-proof}
contains the complete representation argument.

For $q\in\D(A)$ and $K:A\to\D(O\times R)$, write
$E_{qK}:=\E_{qK}[u(O)]$ and, whenever $p_{qK}(o,r)>0$,
\[
 \pi^{qK}_{or}(a):=\frac{q(a)K(o,r\mid a)}{p_{qK}(o,r)},
 \qquad
 \mu_{qK}:=\sum_{o,r}p_{qK}(o,r)\,\delta_{\pi^{qK}_{or}},
\]
where zero-probability terms are omitted.

\subsection{Consequences of the representation}
\label{app:representation-consequences}

\begin{corollary}[No payoff--trace complementarity]\label{cor:matching}
Suppose Axioms \textup{(A1)--(A8)} hold.  Fix $A$ and $q\in\D(A)$.  If channels $K$ and $L$ satisfy
$E_{qK}=E_{qL}$ and $\mu_{qK}=\mu_{qL}$, then $(q,K)\sim(q,L)$.
\end{corollary}

\begin{proof}
The result follows from Theorem~\ref{thm:main} and
$\I_{qK}(A;O,R)=\Sh(q)-\int\Sh(\pi)d\mu_{qK}(\pi)$.
\end{proof}

\begin{proof}[Proof of Proposition~\ref{prop:uniqueness}]
\emph{Range.}
Choose $o_*\in O$ with $u(o_*)=\underline u$.  The range of $V$ is
$[\underline u,\infty)$.  The lower bound is immediate.  Let a constant-$o_*$
channel reveal the realised action with probability $t$ and otherwise report $*$:
\[
 E_t(o_*,a'\mid a)=t\1_{a'=a},
 \qquad
 E_t(o_*,*\mid a)=1-t.
\]
At the uniform lottery on $n$ actions,
$V(q,E_t)=\underline u+\lambda t\log n$ fills
$[\underline u,\underline u+\lambda\log n]$; these intervals are unbounded.

\emph{Ordinal uniqueness.}
Let $W$ be common-scale and define
$\varphi(v):=W(q,K)$ for any pair with $V(q,K)=v$.  This is well defined: equality of the two $V$-values gives
indifference in their two-block environment by Theorem~\ref{thm:main}, and the common-scale property gives equality
of their $W$-values; strict inequalities are preserved in the same way.  Thus
$W=\varphi\circ V$ for a strictly increasing $\varphi$.  The converse follows from
Theorem~\ref{thm:main}, proving part~\textup{(i)}.

\emph{Affine uniqueness.}
The two chain rules give, for every
$P:A\to\D(\{1,\ldots,k\})$ and every
$K_1,\ldots,K_k,L_1,\ldots,L_k:A\to\D(O\times R)$, writing
$m_i:=\sum_aq(a)P(i\mid a)$,
\[
 V\bigl(q,P\triangleright(K_1,\ldots,K_k)\bigr)
 -V\bigl(q,P\triangleright(L_1,\ldots,L_k)\bigr)
 =\sum_{i:m_i>0}m_i
 \bigl[V(q_i,K_i)-V(q_i,L_i)\bigr].
\]
Hence every $\alpha V+\beta$, $\alpha>0$, is common-scale and branch-additive.

Conversely, write a branch-additive common-scale representation as
$W=\varphi\circ V$ and put
\[
 \psi(x):=\varphi(\underline u+x)-\varphi(\underline u),
 \qquad x\ge0.
\]
For $x\in[0,\lambda\log n]$, use the range construction after the first
branch of a public coin with probabilities $\theta$ and $1-\theta$, and the
uninformative constant-$o_*$ channel elsewhere.  Comparing this compound with
the uninformative compound in \eqref{eq:branch-additive-index} gives
\[
 \psi(\theta x)=\theta\psi(x)
 \qquad(0\le\theta\le1).
\]
Taking $x=\lambda\log n$ shows that $\psi$ is linear on $[0,\lambda\log n]$.  These intervals are nested and
unbounded, so $\psi(x)=\alpha x$ on $[0,\infty)$ for one $\alpha>0$.  Therefore
$\varphi(v)=\alpha v+\beta$, proving part~\textup{(ii)}.

\emph{Parameters.}
Any second trace-tempered index $\widetilde V$ is common-scale by block
invariance and branch-additive by the same chain rules.  Part~\textup{(ii)},
restricted first to deterministic payoffs and then to constant-payoff record channels, gives
$\widetilde u=\alpha u+\beta$ and $\widetilde\lambda=\alpha\lambda$.
\end{proof}

\begin{proof}[Proof of Corollary~\ref{cor:signduality}]
Part \textup{(i)} is Theorem~\ref{thm:main}.  For part \textup{(ii)}, reverse
the primitive relation in every environment.  The reversed family satisfies
Axioms~\textup{(A1)}, \textup{(A2)}, \textup{(A5)}, and \textup{(A8)}; Axiom~\textup{(A3)} holds with
its two witnessing outcomes interchanged.  It also satisfies the forward
versions of Axioms~\textup{(A4)}, \textup{(A6)}, and
\textup{(A7)}.  Theorem~\ref{thm:main} therefore represents the reversed
family by
\(
  \E[\widehat u(O)]+\widehat\lambda\I(A;O,R)
\)
with \(\widehat\lambda>0\).  Negating that index represents the original
family, so it has the asserted form with
\(u=-\widehat u\) and \(\lambda=-\widehat\lambda<0\).

For part \textup{(iii)}, collapse the action alphabet of any pair $(q,K)$ to
a singleton.  Axiom~$(\mathrm{A}7^0)$ makes the pair indifferent to the
resulting one-action channel, whose output law is $p_{qK}$.  Axiom
$(\mathrm{A}6^0)$ then erases its record without changing value.  Thus every
pair is indifferent to the one-action channel that delivers precisely its
payoff marginal.  Axiom~\textup{(A5)} makes these reductions coherent across
blocks.  On payoff lotteries, Axiom~\textup{(A8)} applied to an
action-independent public coin gives mixture independence, Axiom
\textup{(A2)} gives continuity, and Axiom~\textup{(A3)} gives
non-triviality.  The mixture-space theorem therefore yields a nonconstant
expected-utility index.  Conversely, expected utility is unaffected by record
or action processing and satisfies the three indifference clauses.  The
positive and negative formulae satisfy their respective clauses by data
processing and the two chain rules.
\end{proof}

\subsection{Independence of the axioms}\label{app:axiom-independence}

\begin{proof}[Proof of Proposition~\ref{prop:independence}]
Fix nonconstant $u$ and write $U(q,K):=\E_{qK}[u(O)]$, $Z:=(O,R)$, and
$I(q,K):=\I_{qK}(A;Z)$.  All scores below use the same formula in block and
compound environments.  Both coordinates are block invariant and processing
monotone; under a common-branch replacement their differences are multiplied
by the reached-branch probability.  Moreover, $U$ distinguishes the material
test and $I$ takes values $\log2$ and zero in the trace test.  For witnesses
built only from $(U,I)$, these facts verify every unmentioned axiom below.

\emph{Failure of Axiom \textup{(A1)}.}
Use the Pareto relation on the two coordinates $(U,I)$:
\[
 (q,K)\succeq(p,L)
 \quad\Longleftrightarrow\quad
 U(q,K)\ge U(p,L)\ \text{and}\ I(q,K)\ge I(p,L).
\]
It is transitive with closed graph, but incomplete whenever payoff and
information rank two alternatives oppositely.

\emph{Failure of Axiom \textup{(A2)}.}
Order alternatives lexicographically by $(U,I)$, with payoff first.  This is
a weak order but is not continuous even within one fixed channel.  Let a
three-action channel reveal the action fully and let only action $2$ pay the high outcome.  Put
\[
 p=(1/4,1/2,1/4),\qquad q=(0.49,0.50,0.01),\qquad
 q_n=(0.49-\varepsilon_n,0.50+\varepsilon_n,0.01),
\]
where $\varepsilon_n\downarrow0$.  The lotteries $p$ and $q$ have the same expected utility but
$\Sh(q)<\Sh(p)$, so $p\succ_Kq$.  Every $q_n$ has strictly higher expected utility than $p$, so
$q_n\succ_Kp$.  Thus the graph in Axiom~\textup{(A2)} is not closed.

\emph{Failure of Axioms \textup{(A3)} and \textup{(A4)}.}
For Axiom~\textup{(A3)}, use $V:=I$: pure lotteries carry no information,
although the trace test has values $\log2$ and zero.  For
Axiom~\textup{(A4)}, use $V:=U$: suitable outcomes have different utility,
but the two trace-test lotteries have the same payoff.

\emph{Failure of Axiom \textup{(A5)}.}
The idea is to add a linear-in-the-prior bonus that vanishes in every canonical
two-block comparison but becomes active in a three-block environment.
Define the weights directly from the literal labels of each finite action alphabet $A$.  If every element of
$A$ is literally a pair $(x,i)$ and the set of outermost second coordinates is
$\{0,\ldots,m-1\}$ for some $m\ge3$, put
\[
 w_A(x,i):=\1_{\{i=0\}}.
\]
Otherwise put $w_A=0$.  In a nested tagged sum this convention reads only the literal outermost coordinate, so
it is a single-valued function of the actual action set and not of its construction history.  For
$\varepsilon>0$, use within each environment
\[
 V_\varepsilon(q,K):=U(q,K)+I(q,K)+\varepsilon\sum_aq(a)w_A(a).
\]
This is a continuous weak order.  The material- and trace-relevance tests can use action sets on which
$w_A=0$, so Axioms~\textup{(A3)} and~\textup{(A4)} retain their usual witnesses.
Every comparison in Axioms~\textup{(A6)--(A8)} is evaluated in a canonical outer two-block environment;
there the outermost tag set is $\{0,1\}$, even when a component alphabet is itself tagged.  Hence $w_A=0$
throughout those comparisons and the score reduces to $U+I$.  For three identical one-action,
constant-payoff, silent channels, however, blocks $0$ and $1$ are indifferent
in their canonical two-block comparison but differ by $\varepsilon$ in the
three-block environment.  Thus Axiom~\textup{(A5)} fails.

\emph{Failure of Axiom \textup{(A6)}.}
For $\varepsilon>0$, use
\begin{equation}
 V(q,K):=U(q,K)+I(q,K)-\varepsilon\Sh_{qK}(Z\mid A).
 \label{eq:no-record-dp}
\end{equation}
Conditional entropy is continuous and block invariant, and mutual information
and conditional entropy both obey the branch chain rule.  Under action processing $A\to B$, the
original-minus-processed difference is
\[
 \bigl[\I(A;Z)-\I(B;Z)\bigr]
 +\varepsilon\bigl[\Sh(Z\mid B)-\Sh(Z\mid A)\bigr]
 =(1+\varepsilon)\I(A;Z\mid B)\ge0,
\]
so Axiom~\textup{(A7)} holds.  In the trace test the two pairs
$(I(q,K),\Sh(Z\mid A))$ are $(\log2,0)$ and $(0,\log2)$.  But with constant
payoff and a fair record independent of $A$, the score is
$-\varepsilon\log2$; deleting that noise raises it to zero, violating
Axiom~\textup{(A6)}.

\emph{Failure of Axiom \textup{(A7)}.}
Let $G(q):=1-\sum_aq(a)^2$ be Gini entropy and define its posterior reduction
\[
 J_G(q,K):=G(q)-\int G(\pi)\,d\mu_{qK}(\pi),
 \qquad V:=U+J_G.
\]
The score $J_G$ is continuous and block invariant.  Concavity of $G$ gives
record processing, and the potential-difference form gives the branch chain
rule; the trace-test values are $1/2$ and zero.
Action processing can nevertheless increase $J_G$.  Take
$q=(1/8,1/8,3/4)$, a constant payoff, and a binary record with
\[
 \Pr(0\mid a_1)=\Pr(0\mid a_2)=0,
 \qquad \Pr(0\mid a_3)=1/4,
\]
and merge $a_1,a_2$.  Direct calculation gives $J_G=9/416$ before the merge
and $J_G=3/104$ afterwards, violating Axiom~\textup{(A7)}.

\emph{Failure of Axiom \textup{(A8)}.}
The idea is that a cap preserves processing monotonicity but makes a branch's
value depend on how much information has already accumulated elsewhere.
Use $V(q,K):=U(q,K)+\min\{I(q,K),1\}$ in bits.  The cap preserves continuity, block
coherence, processing, and both relevance tests.  Let $q$ be uniform on four
actions and let a first-stage record reveal which of two pairs contains the
action.  Keep payoffs constant and one continuation silent; after the other
record, compare a silent continuation with one that reveals the action within its pair.  The
latter is strictly preferred at the branch posterior, but the silent and revealing compounds have
one and three-halves bits, respectively; both are capped at one and are therefore indifferent.
Taking $K$ silent and $L$ revealing in
\eqref{eq:recordwise-sure-thing}, the compound comparison is true by
indifference whereas $(q_1,K)\not\succeq(q_1,L)$; the biconditional fails.
\end{proof}

\subsection{Choice implications}\label{app:choice-proofs}

The next four blocks mirror Section~\ref{sec:implications}.  Write
$Z:=O\times R$; when $K:A\to\D(Z)$ is fixed, put
$K_a:=K(\cdot\mid a)$ and
$\bar u(a):=\E_{K_a}[u(O)]$; its output law under $q$ is $p_{qK}$.

\subsubsection{The trace test}\label{app:trace-test-proof}

\begin{proof}[Proof of Proposition~\ref{prop:signature}]
(i) Under either lottery, each visible pair $(o_0,r_j)$, $j=1,2$, has probability $\tfrac12$: under $q'$ because
each $a_j$ produces $(o_0,r_j)$ for sure, under $q''$ because each of $a_3,a_4$ produces either pair with equal
probability.  Hence $p_{q'K}=p_{q''K}$.  On $\supp(q')$ the rows are degenerate and distinct, so the visible pair
reveals the action and $\I_{q'K}=\Sh(q')=\log2$; on $\supp(q'')$ the rows are identical, so the visible pair is
independent of the action and $\I_{q''K}=0$.  The payoff terms agree because the payoff marginals do, giving the
displayed gap.
(ii) All rows of $K^\circ$ are equal, so $\I_{qK^\circ}=0$ for every $q$; since the payoff is constantly $o_0$,
the representation gives indifference.
\end{proof}

\subsubsection{Deliberate randomisation}
\label{app:randomisation-proofs}

\begin{proposition}[Optimal mixing]\label{prop:binary}
Let $\lambda>0$ and $A=\{a,b\}$ with $K_a\ne K_b$ and
$\Delta:=\bar u(a)-\bar u(b)\ge0$.  For $x\in[0,1]$ let $q_x$ put mass $x$ on $b$, so that
$p_{q_xK}=(1-x)K_a+xK_b=:p_x$, and for $x\in(0,1)$ define
\[
 g(x):=\KL\!\left(K_b\,\middle\|\,p_x\right)-\KL\!\left(K_a\,\middle\|\,p_x\right).
\]
\begin{enumerate}[label=\textup{(\roman*)},leftmargin=*]
\item $g$ is continuous and strictly decreasing, with
$g(0^+)=\KL(K_b\,\|\,K_a)$ and
$g(1^-)=-\KL(K_a\,\|\,K_b)$ as extended-real limits.  The pure lottery $\delta_a$ is optimal if and only if
$\Delta\ge\lambda\KL(K_b\,\|\,K_a)$.  Below this threshold the optimum is unique and interior,
$q^*=q_{x^*}$ with $x^*$ the unique solution of $\Delta=\lambda g(x)$.
\item The optimal mass $x^*(\Delta,\lambda)$ on the inferior action is nonincreasing in $\Delta$ and
nondecreasing in $\lambda$.  While the optimum is interior, it is strictly decreasing in $\Delta$ and,
when $\Delta>0$, strictly increasing in $\lambda$.  At $\Delta=0$ it is the capacity-achieving input
of the two-row channel, independently of $\lambda$.
\end{enumerate}
\end{proposition}

\begin{proof}
Write
\[
 f(x):=V(q_x,K)=\bar u(a)-x\Delta+\lambda I(x),
 \qquad
 I(x):=\Sh(p_x)-(1-x)\Sh(K_a)-x\Sh(K_b).
\]
For $x\in(0,1)$, differentiation gives $I'(x)=g(x)$ and
\[
 g'(x)
 =-c_{\log}\sum_{z\in Z}\frac{(K_b(z)-K_a(z))^2}{p_x(z)}<0,
\]
where $c_{\log}=1/\ln b_0>0$ for the fixed logarithm base $b_0$
and zero summands are omitted.  Thus $f'(x)=-\Delta+\lambda g(x)$ is strictly
decreasing.  As $x\downarrow0$, finite summands converge to those in
$\KL(K_b\|K_a)$, while any $z$ with $K_b(z)>0=K_a(z)$ makes the limit
$+\infty$; exchanging $a$ and $b$ gives the limit at one.  Consequently
$x=0$ is optimal exactly when
$\Delta\ge\lambda\KL(K_b\|K_a)$.  Otherwise $f'(0^+)>0$ and
$f'(1^-)=-\Delta-\lambda\KL(K_a\|K_b)<0$, so the unique optimum is the
root $\Delta=\lambda g(x)$.  At an interior optimum,
$x^*=g^{-1}(\Delta/\lambda)$; since $g^{-1}$ is strictly decreasing, this
gives strict decrease in $\Delta$ and, when $\Delta>0$, strict increase in
$\lambda$, together with the stated weak comparative statics at the boundary.
At $\Delta=0$, $g(x)=0$ is the usual capacity condition and is independent
of $\lambda$.
\end{proof}

\paragraph{General optimality.}
For every $q$, the entropy identity gives
\begin{equation}\label{eq:choice-value}
 \begin{aligned}
 V(q,K)
 &=\lambda\Sh(p_{qK})
   +\sum_aq(a)\bigl[\bar u(a)-\lambda\Sh(K_a)\bigr] \\
 &=\sum_{a\in\supp(q)}q(a)
   \left[\bar u(a)+\lambda
   \KL\!\left(K_a\,\middle\|\,p_{qK}\right)\right].
 \end{aligned}
\end{equation}
Define
\[
 \Phi_a(p):=\bar u(a)+\lambda\KL\!\left(K_a\,\middle\|\,p\right).
\]
For $q_t:=(1-t)q+t\delta_a$, direct differentiation in the extended-real
sense gives
\begin{equation}\label{eq:choice-direction}
 \left.\frac{d}{dt}V(q_t,K)\right|_{t=0^+}
 =\Phi_a(p_{qK})-V(q,K).
\end{equation}

\begin{proof}[Proof of Proposition~\ref{prop:optimal-choice}]
Expected utility is linear in $q$, while mutual information is concave in
the input distribution \cite{CoverThomas2006}; continuity on the compact
simplex gives existence.  At an optimum $q^*$, every vertex direction in
\eqref{eq:choice-direction} is nonpositive, so
$\Phi_a(p_{q^*K})\le V(q^*,K)$ for all $a$.  Equation
\eqref{eq:choice-value} says their $q^*$-weighted average is $V(q^*,K)$;
hence equality holds on $\supp(q^*)$.

Conversely, the proposition's equalities and inequalities, averaged under
$q^*$, give $c=V(q^*,K)$.  The derivative toward any lottery $r$ is
\[
 \sum_a r(a)\bigl[\Phi_a(p_{q^*K})-c\bigr]\le0.
\]
Concavity makes this condition sufficient for global optimality.  At
constant $u$ it is the classical capacity condition.
\end{proof}

\begin{proposition}[Support and pure choice]\label{cor:support-pure}
Let $\lambda>0$ and let
$Z_K:=\{z\in Z:K_a(z)>0\text{ for some }a\in A\}$.
\begin{enumerate}[label=\textup{(\roman*)},leftmargin=*]
\item Every optimal lottery $q^*$ satisfies $p_{q^*K}(z)>0$ for every
$z\in Z_K$.  Consequently, if every action $a$ has a private
consequence---some $z_a$ with $K_a(z_a)>0$ and $K_b(z_a)=0$ for all
$b\ne a$---then every optimal lottery has full support.
\item The pure lottery $\delta_a$ is optimal if and only if
\[
 \bar u(b)+\lambda\KL\!\left(K_b\,\middle\|\,K_a\right)
 \le \bar u(a)
 \qquad\text{for every }b\in A,
\]
and uniquely optimal if these inequalities are strict for $b\ne a$.  In
particular, if some rival reaches a consequence outside the support of
$K_a$, then $\delta_a$ is not optimal, however large its payoff advantage.
\end{enumerate}
\end{proposition}

\begin{proof}
(i) Put $p^*:=p_{q^*K}$.  If $p^*(z)=0<K_a(z)$, then
$\KL(K_a\|p^*)=+\infty$, contradicting either the on-support equality or
the off-support inequality in Proposition~\ref{prop:optimal-choice}.
If action $a$ has a private consequence but $q^*(a)=0$, that consequence
has zero probability under $p^*$, contradicting the first claim.

(ii) At $q=\delta_a$ the output law is $K_a$ and the on-support value is
$\bar u(a)$.  Proposition~\ref{prop:optimal-choice} therefore gives exactly
the displayed inequalities.  If they are strict off $a$, every nonzero
direction from $\delta_a$ has strictly negative directional derivative, so
concavity makes $\delta_a$ the unique maximiser.
If $K_b$ is not supported by $K_a$, then $\KL(K_b\|K_a)=+\infty$, so the
displayed condition fails.
\end{proof}

\begin{proof}[Proof of Corollary~\ref{rem:logit}]
Pairwise-disjoint row supports make the visible consequence identify the
action, so $\I_{qK}(A;Z)=\Sh(q)$; they also give every action a private
consequence, so Proposition~\ref{cor:support-pure}\textup{(i)} gives full
support at the optimum.  On $\supp K_a$,
$p_{qK}(z)=q(a)K_a(z)$.  Hence
$\KL(K_a\|p_{qK})=-\log q(a)$.  The equalities in
Proposition~\ref{prop:optimal-choice} give
$\bar u(a)-\lambda\log q^*(a)=c$.  In natural-log units, solving and
normalising yields the displayed logit formula.  Strict concavity of
$\sum_aq(a)\bar u(a)+\lambda\Sh(q)$ gives uniqueness; changing the
logarithm base merely rescales $\lambda$.
\end{proof}

\begin{proposition}[Optimal marginal and multiplicity]\label{prop:optimal-marginal}
Let $\lambda>0$ and write $C^*(K):=\max_{q\in\D(A)}V(q,K)$.  There is a unique $p^*\in\D(Z)$ such that
$p_{q^*K}=p^*$ for every optimal lottery $q^*$.  With
\[
 \mathcal A^*
 :=\bigl\{a\in A:
 \bar u(a)+\lambda
 \KL\!\left(K_a\,\middle\|\,p^*\right)=C^*(K)
 \bigr\},
\]
the set of optimal lotteries is exactly
\[
 \bigl\{q\in\D(A):
 p_{qK}=p^*,\ \supp(q)\subseteq\mathcal A^*
 \bigr\}.
\]
The optimum is unique whenever the rows $\{K_a:a\in\mathcal A^*\}$ are affinely independent.  If actions are
partitioned into classes with identical rows, $V(q,K)$ depends on $q$ within each class only through the class's
total mass.  The optimal lotteries are exactly the lifts of the optimal lotteries of the collapsed channel, and
each class's mass may be distributed arbitrarily among its members.
\end{proposition}

\begin{proof}[Proof of Proposition~\ref{prop:optimal-marginal}]
If optimal $q_0,q_1$ induced different output laws, the first line of
\eqref{eq:choice-value} would be strictly concave along their segment while
its remaining term is linear, producing a value above $C^*(K)$.  Thus every
optimum induces the same $p^*$.  Proposition~\ref{prop:optimal-choice} and
\eqref{eq:choice-value} then show that every active score equals $C^*(K)$,
so $\supp(q^*)\subseteq\mathcal A^*$.  Conversely, if $p_{qK}=p^*$ and
$\supp(q)\subseteq\mathcal A^*$, then
$V(q,K)=\sum_aq(a)\Phi_a(p^*)=C^*(K)$.
Affine independence makes the representation $p^*=\sum_aq(a)K_a$ with $\supp(q)\subseteq\mathcal A^*$ unique.
Identical rows share $\bar u(a)$ and $\Sh(K_a)$, so both $p_{qK}$ and
\eqref{eq:choice-value} are
unchanged by redistribution within a class; collapsing the classes and lifting aggregate masses proves the last
claim.
\end{proof}

\begin{remark}[Overlapping traces and IIA]\label{rem:overlap-iia-example}
Let three actions have equal expected utility and rows
\[
 K_1=(1,0,0),\qquad K_2=(0,1,0),\qquad K_3=(1/2,0,1/2).
\]
With only the first two actions, the optimal lottery is $(1/2,1/2)$.  Adding the third gives
\[
 q^*=(1/3,4/9,2/9),\qquad p_{q^*K}=(4/9,4/9,1/9).
\]
Each active row has divergence $\log(9/4)$ from $p_{q^*K}$, so
Proposition~\ref{prop:optimal-choice} verifies optimality.  The ratio of the first two choice probabilities falls
from $1$ to $3/4$.  Thus overlapping rows can generate the red-bus/blue-bus pattern even when expected utilities
are equal.
\end{remark}

\subsubsection{The value of being on record}
\label{app:record-value-proofs}

\begin{proof}[Proof of Proposition~\ref{prop:garbling-price}]
The processor leaves the payoff coordinate unchanged, so the expected-utility terms coincide and the difference is
$\lambda[\I_{qK}(A;O,R)-\I_{qL}(A;O,R')]$.  Write $Z:=(O,R)$ and $Z':=(O,R')$.  By construction $A\to Z\to Z'$ is
a Markov chain, so expanding $\I(A;Z,Z')$ in both orders gives
\[
 \I(A;Z)-\I(A;Z')=\I(A;Z\mid Z')=\I(A;R\mid O,R'),
\]
the last equality because $O$ is part of $Z'$.  Conditional mutual information is nonnegative and vanishes exactly
when $A$ and $(O,R)$ are conditionally independent given $(O,R')$ \cite{CoverThomas2006}.  This is the zero-price
condition stated in the proposition.  Under a garbling, it is also equivalent to equality of the posterior laws,
$\mu_{qK}=\mu_{qL}$.  Indeed, the Markov property gives
$\Pr(A=\cdot\mid Z,Z')=\pi^{qK}_{Z}$ almost surely, while conditional independence gives
$\Pr(A=\cdot\mid Z,Z')=\pi^{qL}_{Z'}$; the two random posteriors then coincide almost surely.  Conversely, by the
identity in the proof of Corollary~\ref{cor:matching}, equality of the posterior laws gives equality of the
information terms, and the payoff terms already agree.
\end{proof}

\begin{corollary}[The recorded coin]\label{cor:recorded-coin}
Let $K,L:A\to\D(O\times R)$ and $t\in(0,1)$.  Let $tK\oplus(1-t)L$ be the public mixture that draws an
action-independent coin $Y$ with $\Pr(Y=0)=t$, applies $K$ if $Y=0$ and $L$ otherwise, and includes the coin in
the record, so the full public-mixture record is $(Y,R)$; let $M_t:=tK+(1-t)L$ be the rowwise mixture, which hides it.  Then, for every $q\in\D(A)$,
\begin{enumerate}[label=\textup{(\roman*)},leftmargin=*]
\item $V\bigl(q,\,tK\oplus(1-t)L\bigr)=t\,V(q,K)+(1-t)\,V(q,L)$;
\item $V\bigl(q,\,tK\oplus(1-t)L\bigr)-V(q,M_t)=\lambda\,\I(A;Y\mid O,R)\ge0$, with equality if and only if $A$
and $Y$ are conditionally independent given $(O,R)$.
\end{enumerate}
\end{corollary}

\begin{proof}[Proof of Corollary~\ref{cor:recorded-coin}]
Part \textup{(i)} is the chain rule with an action-independent first stage:
$\I(A;Y)=0$, and each reached branch reproduces the joint law of the
corresponding constituent.  Part \textup{(ii)} follows from
Proposition~\ref{prop:garbling-price}: deleting the coin coordinate is a
deterministic payoff-preserving record processor carrying the public mixture
to $M_t$.
\end{proof}

\begin{proposition}[Uniform dominance is more-capable, not garbling]
\label{prop:blackwell}\label{def:channel-orders}
Let $K:A\to\D(O\times R)$ and $L:A\to\D(O\times R')$ have identical
payoff kernels.  Write $K\succeq_gL$ if $L=KT$ for some record processor
$T$, and write $K\succeq_cL$ if
$\I_{qK}(A;O,R)\ge\I_{qL}(A;O,R')$ for every $q\in\D(A)$.  The latter is
the \emph{more-capable} order \cite{KornerMarton1977}.
\begin{enumerate}[label=\textup{(\roman*)},leftmargin=*]
\item $K\succeq_cL$ if and only if $V(q,K)\ge V(q,L)$ for every $q\in\D(A)$.
\item $K\succeq_gL\Rightarrow K\succeq_cL$, but the converse fails.  There are channels for which
$K\succeq_cL$ although $L\ne KT$ for every record processor $T$.
\end{enumerate}
\end{proposition}

\begin{proof}[Proof of Proposition~\ref{prop:blackwell}]
(i) Because the payoff kernels coincide,
$V(q,K)-V(q,L)=\lambda[\I_{qK}-\I_{qL}]$ with $\lambda>0$.  Uniform value dominance and
$K\succeq_cL$ are therefore the same statement.

(ii) If $L=KT$, Proposition~\ref{prop:garbling-price} gives $\I_{qK}\ge\I_{qL}$ at every $q$, and hence
$K\succeq_gL\Rightarrow K\succeq_cL$.  For failure of the converse, let $A=\{0,1\}$, fix $\bar o\in O$, and let
both channels pay $\bar o$ surely.  Let $K$ carry the binary
erasure record with erasure probability $\varepsilon=\tfrac12$ and $L$ the binary symmetric record with crossover
$p=\tfrac15$:
\[
 K(\bar o,a\mid a)=1-\varepsilon,\quad K(\bar o,\mathrm e\mid a)=\varepsilon,
 \qquad
 L(\bar o,a\mid a)=1-p,\quad L(\bar o,1-a\mid a)=p.
\]
The payoff kernels agree row by row, so it suffices to compare informations, and the comparison is invariant to the
base; compute in bits with $h$ the binary entropy.  For $q=(1-t,t)$, grouping records gives
\[
 \I_{qK}=(1-\varepsilon)\,h(t),
 \qquad
 \I_{qL}=h(p\ast t)-h(p),
 \qquad p\ast t:=p(1-t)+(1-p)t.
\]
By Mrs Gerber's Lemma \cite{WynerZiv1973}, $v\mapsto h\bigl(p\ast h^{-1}(v)\bigr)$ is convex on $[0,1]$, with
$h^{-1}$ valued in $[0,\tfrac12]$; since $h(p\ast t)$ is invariant under $t\mapsto1-t$, take $t\le\tfrac12$ and
$v=h(t)$.  The chord between $v=0$ and $v=1$ gives
\[
 h(p\ast t)\le(1-v)\,h(p)+v,
 \qquad\text{hence}\qquad
 \I_{qL}\le\bigl(1-h(p)\bigr)h(t)\le(1-\varepsilon)\,h(t)=\I_{qK},
\]
the last step because $\varepsilon=\tfrac12\le h(\tfrac15)$.  Thus $K\succeq_cL$.  Now suppose $L=KT$ for a
processor $T$, and write $t_r:=T(0\mid \bar o,r)$ for $r\in\{0,1,\mathrm e\}$.  Matching the probability of the
processed record $0$ in the two rows,
\[
 (1-\varepsilon)t_0+\varepsilon t_{\mathrm e}=1-p,
 \qquad
 (1-\varepsilon)t_1+\varepsilon t_{\mathrm e}=p,
\]
and subtracting gives $t_0-t_1=(1-2p)/(1-\varepsilon)=\tfrac65>1$, impossible.  Compositions of record processors
are record processors, so no chain of garblings produces $L$ from $K$ either.
\end{proof}

\subsubsection{Measuring the trace weight}
\label{app:measurement-proofs}

\begin{proposition}[Single information price]
\label{prop:single-price}\label{lem:info-payoff-plane}
If $(q,K)\sim(p,L)$ and $\I_{qK}\ne\I_{pL}$, then
\[
 \frac{E_{pL}-E_{qK}}{\I_{qK}-\I_{pL}}=\lambda,
\]
and $\lambda$ is the only number satisfying this identity at every such
indifference.  More generally, writing
\[
 \phi(q,K):=\bigl(\I_{qK}(A;O,R),E_{qK}\bigr),\qquad
 \underline u:=\min_o u(o),\quad \overline u:=\max_o u(o),
\]
the image of $\phi$ over all finite pairs is
$[0,\infty)\times[\underline u,\overline u]$.  In these coordinates the
block-comparison indifference classes are the nonempty intersections of the
lines $E=c-\lambda\I$ with that strip.
\end{proposition}

\begin{proof}
The ratio identity rearranges
$E_{qK}+\lambda\I_{qK}=E_{pL}+\lambda\I_{pL}$.
For the image claim, only the reverse inclusion needs proof.  Given
$(t,w)$ in the strip, choose $n$ with $\log n\ge t$ and, by continuity of
entropy between a point mass and the uniform lottery, choose $q_t$ on
$\{1,\ldots,n\}$ with $\Sh(q_t)=t$.  Choose $o^+,o^-\in O$ with
$u(o^+)=\overline u$ and $u(o^-)=\underline u$.  Let a channel reveal the
action in its record and, independently of the action, pay $o^+$ with
probability $\beta$ and $o^-$ otherwise, where
$\beta u(o^+)+(1-\beta)u(o^-)=w$.  It realises
$\phi(q_t,K)=(t,w)$, and Theorem~\ref{thm:main} gives the stated lines.

For uniqueness, if $\lambda'\ne\lambda$, put
$t_0:=\tfrac12\min\{1,(\overline u-\underline u)/\lambda\}>0$.
The image claim supplies pairs with coordinates
$(t_0,\overline u-\lambda t_0)$ and $(0,\overline u)$.  They are indifferent,
whereas their $E+\lambda'\I$ values differ by
$(\lambda'-\lambda)t_0\ne0$.
\end{proof}

\begin{proposition}[Willingness to pay for trace]
\label{prop:wtp}\label{def:dilution}
Let $o^+,o^-$ attain the maximum and minimum of $u$, and put
\[
 \Delta u:=u(o^+)-u(o^-)>0,
 \qquad
 w(\beta):=\beta u(o^+)+(1-\beta)u(o^-).
\]
For a channel $H:C\to\D(O\times T)$, zero-pad it to
$T^\ast:=T\sqcup\{\ast\}$.  For $\varepsilon\in[0,1)$ and
$\beta\in[0,1]$, let $P_\varepsilon:C\to\D(\{1,2\})$ satisfy
$P_\varepsilon(2\mid c)=\varepsilon$ for every $c$, and let $B_\beta$ have
constant rows
\[
 B_\beta(o^+,\ast\mid c)=\beta,
 \qquad B_\beta(o^-,\ast\mid c)=1-\beta.
\]
Write
$D_{\varepsilon,\beta}H:=P_\varepsilon\triangleright(H,B_\beta)$.
Fix $q\in\D(A)$, $p\in\D(B)$, and channels $K:A\to\D(O\times R)$ and $L:B\to\D(O\times S)$.
\begin{enumerate}[label=\textup{(\roman*)},leftmargin=*]
\item For every $\varepsilon\in[0,1)$ and $\beta\in[0,1]$,
\[
 V\bigl(q,\,D_{\varepsilon,\beta}K\bigr)=(1-\varepsilon)\,V(q,K)+\varepsilon\,w(\beta).
\]
\item If $\varepsilon\in(0,1)$ and $\beta,\beta'\in[0,1]$ satisfy
$\bigl(q,D_{\varepsilon,\beta}K\bigr)\sim\bigl(p,D_{\varepsilon,\beta'}L\bigr)$, then
\[
 \frac{\varepsilon}{1-\varepsilon}\bigl[w(\beta)-w(\beta')\bigr]
 =V(p,L)-V(q,K).
\]
When the material parts agree, $E_{qK}=E_{pL}$, the left-hand side equals
$\lambda\bigl[\I_{pL}(B;O,S)-\I_{qK}(A;O,R)\bigr]$: information is priced linearly at rate $\lambda$.
\item Such a pair $(\beta,\beta')$ exists if and only if
$\lvert V(p,L)-V(q,K)\rvert\le\frac{\varepsilon}{1-\varepsilon}\Delta u$; a boundary response therefore brackets
the value difference instead of measuring it.
\end{enumerate}
\end{proposition}

\begin{proof}
(i) The material part of the compound is $(1-\varepsilon)\E_{qK}[u(O)]+\varepsilon w(\beta)$.  For the trace part,
the coin is action-independent, so both branch posteriors equal $q$ and the chain rule for the compound gives an
information term of $(1-\varepsilon)\I_{qK}(A;O,R)$, the side-bet branch contributing zero.
(ii) By the moreover clause of Theorem~\ref{thm:main}, the indifference holds if and only if
$(1-\varepsilon)V(q,K)+\varepsilon w(\beta)=(1-\varepsilon)V(p,L)+\varepsilon w(\beta')$; rearrange, and
substitute the representation when the material parts cancel.
(iii) $w(\beta)-w(\beta')$ is continuous on $[0,1]^2$ with range $[-\Delta u,\Delta u]$; apply the intermediate
value theorem to the equation in \textup{(ii)}.
\end{proof}

\begin{proof}[Proof of Proposition~\ref{prop:elicit}]
Write $K^{\mathrm{id}}$ and $K^{0}$ for the revealing and the silent channel at the constant payoff, so that
$V(q,K^{\mathrm{id}})-V(q,K^{0})=\lambda\Sh(q)$.  By Proposition~\ref{prop:wtp}\textup{(i)}, the two diluted
values are $(1-\varepsilon)V(q,K^{\mathrm{id}})+\varepsilon w(0)$ and
$(1-\varepsilon)V(q,K^{0})+\varepsilon w(\beta)$, and under the normalisation $w(\beta)=\beta$.  Indifference
therefore holds if and only if $\varepsilon\beta=(1-\varepsilon)\lambda\Sh(q)$.  The left side is continuous and
strictly increasing in $\beta$ with range $[0,\varepsilon]$, so a solution exists, and is then unique, exactly
when $\lambda\Sh(q)\le\varepsilon/(1-\varepsilon)$; when this fails, the revealing side's value exceeds the
silent side's at every $\beta$.
\end{proof}

\begin{remark}[Implementation and controls]\label{rem:elicitation}
Normalise $u(o^+)=1$ and $u(o^-)=0$.  All comparisons below are primitive
block comparisons.  For each $o\in O$, elicit $\beta_o$ from
$(\delta_\ast,K_o)\sim(\delta_\ast,B_{\beta_o})$, where $K_o$ pays $o$
surely; both one-action channels carry zero information, so
$u(o)=\beta_o$.  Next take the uniform $q$ on $n\ge2$ actions and the
constant-payoff revealing and silent channels $K_o^{\mathrm{id}}$ and
$K_o^0$, and elicit $\beta$ from
\[
 (q,D_{\varepsilon,\beta}K_o^0)
 \sim(q,D_{\varepsilon,0}K_o^{\mathrm{id}}),
 \qquad
 \lambda=\frac{\varepsilon}{1-\varepsilon}\frac{\beta}{\log n}.
\]
If even $\beta=1$ does not produce indifference, the response gives the lower bound
$\lambda>\varepsilon/((1-\varepsilon)\log n)$, and $\varepsilon$ is raised.

Three controls are essential.  \emph{Orientation:} revelation must beat
silence at every constant payoff.  \emph{Garbling:} in a separate comparison of the undiluted,
equal-payoff channels $K_o^{\mathrm{id}}$ and $K_o^0$, total record
erasure must make the two blocks indifferent; otherwise some display feature
still traces the action.  \emph{Transport:} changing the prior, payoff, or
action-set size must return the same $\lambda$ by
Proposition~\ref{prop:single-price}; a decline with $\log n$ instead signals
the capped-information criterion.
\end{remark}
